% Paquetes del informe de correcciones — compatibles con pdflatex
\usepackage[T1]{fontenc}
\usepackage[utf8]{inputenc}
\usepackage[spanish,es-tabla]{babel}
\usepackage[a4paper,margin=2.5cm]{geometry}
\usepackage{graphicx}
\usepackage{booktabs}
\usepackage{tabularx}
\usepackage{longtable}
\usepackage{multirow}
\usepackage{array}
\usepackage{siunitx}
\usepackage{caption}
\usepackage{subcaption}
\usepackage{hyperref}
\usepackage{cleveref}
\usepackage{amsmath}
\usepackage{xcolor}
\usepackage{enumitem}
\usepackage{float}
\usepackage{listings}

\sisetup{
  round-mode = places,
  round-precision = 4,
  detect-all
}

\hypersetup{
  colorlinks = true,
  linkcolor  = blue!50!black,
  citecolor  = green!40!black,
  urlcolor   = blue!60!black
}

\captionsetup{font=small, labelfont=bf}
\graphicspath{{figures/}}

\lstset{
  basicstyle=\ttfamily\small,
  breaklines=true,
  frame=single,
  backgroundcolor=\color{gray!5}
}
